% Pista pro-24j-q4 · Probabilitat
% Procedència: PAU Catalunya 2024 · Convocatòria de juny · Sèrie 1 · Qüestió 4
\documentclass[11pt,a4paper]{article}
\usepackage{pista}
\pistainfo{Catalunya 2024}{Convocatòria de juny}{Sèrie 1}

\begin{document}

\section*{\textcolor{blauPista}{Qüestió 4}}

\noindent L'Andreu posa les nou boles que es mostren a continuació dins d'una bossa.

\noindent\textit{(A la figura: nou boles amb les lletres B, A, Y, E, S, F, A, N, S.)}

\subsection*{\textit{a)} \normalfont A continuació, treu de la bossa dues boles a l'atzar, una darrere l'altra i sense reemplaçament (és a dir, no retorna a la bossa la primera bola abans de treure la segona).}

\noindent --- Calculeu la probabilitat que la primera bola sigui una A o una E. \hspace*{\fill}\textnormal{[0,5~punts]}

\pista{Compta quantes de les nou boles tenen una A i quantes tenen una E. Vigila, perquè alguna lletra surt més d'una vegada. Llavors, aplica la regla de Laplace:
\[
P = \frac{\text{casos favorables}}{\text{casos possibles}}
\]}

\noindent --- Calculeu la probabilitat que les dues boles siguin diferents. \hspace*{\fill}\textnormal{[0,75~punts]}

\pista{És més fàcil calcular la probabilitat de l'esdeveniment contrari:
\[
P(\text{diferents}) = 1 - P(\text{iguals})
\]
Per calcular $P(\text{iguals})$, comença per les A. La probabilitat de treure dues A seguides és $\frac{2}{9} \cdot \frac{1}{8}$, perquè, després de treure la primera A, a la bossa només queda una A entre 8 boles.

\medskip
Després, fes el mateix amb les altres lletres que es repeteixen i, finalment, suma-ho tot.}

\subsection*{\textit{b)} \normalfont L'Andreu torna a posar totes les boles a la bossa i en treu cinc a l'atzar, una darrere l'altra, però ara amb reemplaçament (és a dir, ara sí que retorna a la bossa cada bola extreta abans d'agafar la següent).}

\noindent --- Calculeu la probabilitat que no hagi tret cap A. \hspace*{\fill}\textnormal{[0,5~punts]}

\pista{Com que cada bola es torna a la bossa, a cada extracció hi ha les mateixes nou boles. Per tant, la probabilitat de no treure cap A és la mateixa a les cinc extraccions.

\medskip
La probabilitat de «cap A» és el producte d'aquestes cinc probabilitats, és a dir, aquesta probabilitat elevada a 5.}

\noindent --- Calculeu la probabilitat que hagi tret almenys dues A. \hspace*{\fill}\textnormal{[0,75~punts]}

\pista{Sigui $X$ el nombre de boles amb una A en les cinc extraccions. Com que les extraccions són independents, $X$ segueix una distribució binomial: $X \sim B(5, p)$, on $p$ és la probabilitat de treure una A en una extracció.

\medskip
«Almenys dues A» és l'esdeveniment contrari de «cap A o exactament una A». Per tant:
\[
P(X \ge 2) = 1 - \bigl[P(X = 0) + P(X = 1)\bigr]
\]
Aquest càlcul suposa molta menys feina que si haguessis de sumar $P(X = 2) + P(X = 3) + P(X = 4) + P(X = 5)$.

\medskip
Recorda la fórmula de la binomial:
\[
P(X = k) = \binom{n}{k}\, p^{k}\, (1 - p)^{n - k}
\]}

\end{document}
